\documentclass[11pt]{article}
\usepackage[a4paper,margin=26mm]{geometry}
\usepackage{amsmath,amssymb,amsthm,mathrsfs,hyperref}
\setlength{\emergencystretch}{3em}
\newtheorem{theorem}{Theorem}
\newtheorem{lemma}[theorem]{Lemma}
\newtheorem{proposition}[theorem]{Proposition}
\newcommand{\MA}{\operatorname{MA}}
\newcommand{\EE}{\mathcal E^1_\theta}
\newcommand{\vol}{\operatorname{vol}}
\title{Strict Borel locality from bounded moving maxima\protect\\
and fixed continuous ramps}
\author{Unreviewed conditional research, continuation 201}
\date{6 October 2026}
\begin{document}
\maketitle
\begin{abstract}
For a big rational bundle, we give a conditional route from its finite-list
complete moving systems to strict locality for every original finite-energy
pair. A moving maximum is identified by an energy squeeze against canonical
majorants. A fixed base twist gives uniform algebraic collars and fixed-test
quadratic bounds, so the squeeze also identifies its raw measure.
Finite nef locality follows from finite ample locality. For unbounded pairs,
fixed continuous ramps select the winning region; uniform ordered
finite-energy errors control the incorrect selections before either limit.
This proves a strict Borel comparison inequality without assuming that the
winning region is open. A separate convex-series argument proves nullness
of poles of every genuine plurisubharmonic function. The genuine scalar and
moving foundations, and the absolute midpoint modulus used for the ordered
errors, remain explicit unproved antecedents. This paper does not identify
all unbounded mixed approximations or all negligible sets and does not
solve original P16.5.
\end{abstract}

\section{Original domain and explicit conditional inputs}
Let $X$ be any smooth integral projective $n$-fold, $n\geq1$, over any
algebraically closed trivially valued characteristic-zero field.
Let $L$ be any actual big rational line bundle, $\theta=c_1(L)$,
$K=X^{\mathrm{an}}$, and $V=\vol(L)>0$.
The genuine cone is the all-point single decreasing FS-net cone on
varying actual bundles with numerical classes tending to $\theta$.
No supported semivaluations are removed. Normalize at the generic trivial
valuation $v_0$ and put $R=P_\theta(0)$ and $B(u)=-E(u)$.
Choose an actual Kodaira identity $L=aA+E_0$, $A$ ample, $E_0$ effective;
write $Z=E_0^{\mathrm{an}}$ and $U=K\setminus Z$.
Clear rational powers before using any sections or ideals.

All results below are conditional on the following whole package.
It is mathematical proof data, not an accepted theorem.
\begin{enumerate}
\item[(S)] The literal genuine cone is convex, finite-max and decreasing
closed at finite generic endpoints and admits genuine all-point envelopes.
Its normalized space is compact in the divisorial-coordinate topology.
The whole outer full-section energy is
\[
 E(u)=\inf_{f\in C(K),\,f\geq u}\mathscr E(f),
 \qquad E(P_\theta(f))=\mathscr E(f).
\]
It is concave, monotone, weakly upper semicontinuous, decreasing-net
continuous and $V$-translation equivariant, with $E(R)=0$.
Every normalized $u\in\EE$ has nonpositive rational original model
majorants $f_i\downarrow u$, with $f_i(v_0)=0$ and
$P_\theta(f_i)\downarrow u$ with decreasing energies.
Perturbed canonical envelopes have the same decreasing continuity.
Every normalized genuine function is at most $R\leq0$.
\item[(M)] A target-independent map $M:\EE\to\mathcal M^+(K)$ has
raw mass $V$, is constant invariant, agrees at model envelopes with
the finite divisorial derivative of $\mathscr E$, has weak canonical
decreasing convergence, and gives $Z$ zero mass.
For every finite list of model obstacles, the complete moving systems
have common actual data $p_m^*L=A_m+T_m$, with $mA_m$ generated
and $mT_m$ effective. Their physical stages and bounded rational
relatively generated model quotients satisfy
\[
 \phi_{m,f}=q_{T_m}+\psi_{m,f},\qquad
 \phi_{m,0}=q_{T_m},\qquad A_m^n\longrightarrow V .
\]
The physical stage is at most $P_\theta(f)$ on all lifts.
The auxiliary energy tends to $E(P_\theta(f))$ and its pushed raw
measure tends weakly to $M(P_\theta(f))$; its generic value tends to zero.
The full auxiliary section trace is its finite intersection energy.
Multiplication by the actual $T_m$ section injects the full auxiliary
section space into the full original one in cleared degrees.
The quotients converge uniformly on compact subsets of $U$ to
$P_\theta(f)-R$. For a nonpositive bounded model obstacle they lie
between its minimum, with a vanishing floor error, and zero.
\item[(A)] Every genuine relative finite-energy test is integrable at
every output. With $\alpha=2^{-n}$ and
$\delta(u,v)=E((u+v)/2)-(E(u)+E(v))/2$, the absolute midpoint estimate is
\[
 \int_U|u-v|\,dM(w)
 \leq C_n\delta(u,v)^\alpha B(w)^{1/2}
       \max\{B(u),B(v),B(w)\}^{1/2-\alpha}.                 \tag{1}
\]
When $n=1$ the last exponent-zero factor is omitted.
The uniform finite-model Taylor estimate passes as a whole fixed
increment to every genuine base and gives both original continuous
derivatives of $E(P_\theta(u+th))$ equal to $\int h\,dM(u)$.
\item[(C)] Genuine relative contraction satisfies
$B((1-r)R+ru)\leq C_n r^2B(u)$ for $0<r<1$.
\end{enumerate}
These are respectively the explicitly unproved scalar and moving
package and the conditional conclusions reconstructed in the completed
round159 paper, Sections 5--10. Continuation200 gives another unreviewed
conditional construction. Neither body is a caller grant.
Here the only accepted NA input is G0 in its actual \emph{ample} domain,
used below only for finite bounded ample locality and original model
testing. Ordinary resolution, nef intersection theory, algebraic positive
products and volume differentiation, and compact Hausdorff Radon facts
are distinguished background.

For clarity, an intrinsic finite-dual-energy target is a positive Radon
measure of raw mass $V$, missing $\{R=-\infty\}$, integrating every
$R-u$, $u\in\EE$, and with
$\sup_{u(v_0)=0}(E(u)+\int(R-u)d\nu)<\infty$.
The pole argument at the end applies to such targets too.

\section{Uniform collars for bounded auxiliary models}
Fix a positive integer $C$ and auxiliary models $\xi$ satisfying
$-C\leq\xi\leq0$ on the generated $A_m$.
Their defining lattice may be taken with all time weights in $[-mC,0]$.
After clearing degrees and resolving its base ideal, a single fixed
integer $c>C+1$ makes
\[
 \Gamma_\xi(c)=p^*L-T_m+D_\xi+cF
\]
globally generated and hence nef on its compactified model.
Here $F$ is the full fiber, with every central multiplicity.
Indeed on the zero chart the weighted section system generates the
resolved lattice; on the infinity chart it is the generated generic
bundle. Twisting by $\mathcal O_{\mathbb P^1}(mc)$ supplies all
monomials between the two charts, since each interval of exponents
has length at least $m(c-C)>0$. This constructs a generating system
on both charts, not merely relative nefness.
The same argument applies to the maximum of two such scalars:
take the union of their weighted generating systems and principalize
its actual ideal. All weights remain in $[-mC,0]$.
In effective order, on every common refinement,
\[
 0\leq\Gamma_\xi(c)\leq\gamma:=p^*L+cF.                  \tag{2}
\]
The right class is big on the fixed product determination; the
horizontal $T_m$ and $-D_\xi$ account for the difference.
For each fixed $m$, this global twist also identifies the full
auxiliary section-volume intercept with $e_{A_m}(\xi)$:
large twists are nef, their volume is their ordinary top intersection,
and its affine fiber term is $cA_m^n$. Thus this finite trace equality
applies to the new maximum as well as to the supplied endpoint models.

\begin{lemma}[Algebraic collars and fixed quadratic tests]\label{collars}
For any fixed actual effective Cartier divisor $S$ on $X$, there is
$C_{S,c}$ independent of $m$ and of these auxiliary models such that
\[
 \int v(S)\,d\MA_{A_m,\mathrm{raw}}(\xi)\leq C_{S,c}.     \tag{3}
\]
After pushforward every weak limit gives $S^{\mathrm{an}}$ zero mass.
For any fixed original rational model test $h=f_H$, its quadratic form
in any $n-1$ of these auxiliary slots is bounded by a constant
$C_{h,c}$ independent of their refinements and $m$.
\end{lemma}
\begin{proof}
On a common SNC determination write $F=\sum_Eb_EE$.
The vertical part of the pulled $S$ has coefficient $b_Ev_E(S)$.
The raw measure coefficient is $b_E(\Gamma_\xi(c)|_E)^n$:
adding $cF$ does not change a central restriction.
The horizontal part of the pulled effective $S$ is effective.
Nef positivity therefore gives
\[
 \int v(S)\,d\MA(\xi)
 \leq \Gamma_\xi(c)^n\,p^*S.
\]
Choose one fixed ample $T$ on a smooth determination and $d$ with
$dT-p^*S$ nef. Algebraic positive-product order and (2) imply
\[
 \Gamma_\xi(c)^n p^*S
 \leq d\Gamma_\xi(c)^n T
 \leq d\langle\gamma^n\rangle T .
\]
Only fixed nef tests are paired with effective-order inequalities.
This proves (3), also for nef classes on the boundary by nef continuity.
The original capped functions $\min(v(S),N)$ are continuous on $K$.
Weak convergence and (3) bound each limiting cap integral by $C_{S,c}$.
Thus the mass of $S^{\mathrm{an}}$, where every cap is $N$, is at most
$C_{S,c}/N$, and is zero. Before any limit the same caps bound the
mass of the collar $\{v(S)\geq N\}$ by $C_{S,c}/N$.

For the quadratic claim choose a single determination carrying $H$.
Write $H=H_1-H_2$ with two \emph{fixed} globally nef Cartier classes,
for example $H_1=aT+H$ and $H_2=aT$ for sufficiently large $a$.
For $\Lambda$ a product of $n-1$ auxiliary slots,
\[
 |-H^2\Lambda|
 \leq\sum_{k,l=1}^2 H_kH_l\Lambda
 \leq\sum_{k,l=1}^2\langle\gamma^{n-1}\rangle H_kH_l .
                                                               \tag{4}
\]
The slots can be replaced by their twists $\Gamma_\xi(c)$ because
every term containing $F$ in a vertical quadratic intersection is zero.
Each term in the middle expression is nonnegative; positive-product
order is used only against the fixed nef products $H_kH_l$.
All choices precede the moving model. For $n=1$ the product is empty
and (4) is an ordinary fixed intersection bound.
\end{proof}

\section{Identifying the bounded moving maximum}
Fix normalized model envelopes $a=P_\theta(f)$ and $b=P_\theta(g)$
of nonpositive model obstacles. Let
\[
 w=\max(a,b),\qquad
 \rho_m=\max(\psi_{m,f},\psi_{m,g}).
\]
The maximum is an actual relatively generated scalar on the same $A_m$.
On the whole pulled carrier its physical stage is
$q_{T_m}+\rho_m=\max(\phi_{m,f},\phi_{m,g})\leq w$.
Choose one integer $C$ so that all endpoint quotient weights and $\rho_m$
lie in $[-C,0]$ and $R-C\leq w\leq R$.
No equality $w=P_\theta(\max(f,g))$ is assumed.

\begin{lemma}[Trace inequality]\label{trace}
For every fixed continuous original majorant $h\geq w$,
\[
 \mathscr E(h)\geq e_{A_m}(\rho_m)
                    +(\min_Kh)(V-A_m^n).               \tag{5}
\]
Consequently $E(w)\geq\limsup_m e_{A_m}(\rho_m)$.
\end{lemma}
\begin{proof}
At fixed $m$ multiply the full auxiliary section system in degree
$\ell A_m$ by the actual $\ell T_m$ section.
Its $h$ filtration weights are at least the weights of $\rho_m$,
because its physical scalar is at most $h$ at every lift.
At supported zero branches the original section expression is
$+\infty$ and cannot lower the weight infimum.
Extend an adapted basis of this injected subspace to a full
original filtered basis. Every remaining weight is at least
$\ell\min_Kh$. First take the degree $\ell$ to infinity at this
fixed $m,h$; the bounded generated-nef trace gives (5).
The floor error divided by $\ell$ tends to zero.
Only after that limit let $m$ tend cofinally to infinity, still with
the same finite $h$. The retained, possibly negative,
$(\min h)(V-A_m^n)$ then vanishes.
Finally take the outer infimum over $h$. No physical stage is
assigned membership in the original $\EE$.
\end{proof}

\begin{theorem}[Bounded maximum energy and measure]\label{maxlimit}
Under (S)--(M),
\[
 e_{A_m}(\rho_m)\longrightarrow E(w),\qquad
 p_{m*}\MA_{A_m,\mathrm{raw}}(\rho_m)\longrightarrow M(w)
                                                        \tag{6}
\]
weakly on the original $K$.
\end{theorem}
\begin{proof}
Take canonical original model majorants $h_j\downarrow w$ and
$w_j=P_\theta(h_j)\downarrow w$.
For fixed $j$ put $\sigma_{m,j}=\psi_{m,h_j}$.
The full filtration order gives $\sigma_{m,j}\geq\rho_m$.
Moreover \emph{the same} $C$ bounds all these $\sigma_{m,j}$:
$h_j\geq w\geq R-C\geq\phi_{m,0}-C$.
Every full original section has zero-obstacle weight zero and its
scalar is at most $\phi_{m,0}$. Thus every $h_j$ weight is at least
$-mC$; its generic upper bound is zero.
Integer $C$ removes a floor error. No bound for $\|h_j\|_\infty$
uniform in $j$ is needed.

The measures of $\rho_m$ have converging raw masses $A_m^n\to V$.
Their energies are bounded above and below: monotonicity bounds them
below by the energy of either endpoint, and all normalized quotient
energies are at most zero. Take any joint convergent subnet with
measure limit $\nu$ and energy limit $e_*$. Lemma~\ref{collars},
with a cleared $E_0$, gives $\nu(Z)=0$.
Finite concavity at the lower potential gives, for fixed $j$,
\[
 0\leq e_{A_m}(\sigma_{m,j})-e_{A_m}(\rho_m)
       \leq\int(\sigma_{m,j}-\rho_m)\,d\MA_{A_m}(\rho_m).
                                                               \tag{7}
\]
On every compact subset of $U$ the changing test tends uniformly
to $w_j-w$. It is between zero and $C$.
Fix a continuous collar in $U$ before taking the subnet; its omitted
mass can be made arbitrarily small using $\nu(Z)=0$ and weak
convergence. This passes the right integral of (7), so
\[
 E(w_j)-e_*\leq\int_U(w_j-w)\,d\nu .
\]
The continuous bounded differences decrease to zero on $U$.
Compact Dini convergence on a fixed collar, followed by its omitted
Radon mass, gives convergence of the last integral to zero even for
the directed $j$ net. Decreasing energy continuity gives
$E(w)\leq e_*$. Lemma~\ref{trace} gives the reverse inequality.
Every joint cluster therefore has energy $E(w)$.

To identify its measure, retain one fixed rational original model
test $h=f_H$. On a common refinement the finite ordered energy gap
$\Delta_{m,j}=e(\sigma_{m,j})-e(\rho_m)$ controls
\[
 I_{m,j}:=\int(\sigma_{m,j}-\rho_m)
                   \,d(\MA(\rho_m)-\MA(\sigma_{m,j}))
        \leq(n+1)\Delta_{m,j}.                           \tag{8}
\]
Indeed the finite energy telescoping sum has $n+1$ nonnegative
terms; its first term is $\int(\sigma-\rho)d\MA(\rho)$.
Subtracting the nonnegative integral at $\sigma$ gives (8).
The same telescoping expresses $I_{m,j}$ as the sum of the
nonnegative mixed quadratic forms of $\sigma-\rho$.
Cauchy--Schwarz in those forms, and the fixed bound (4), therefore give
\[
 \left|\int h\,d(\MA(\rho_m)-\MA(\sigma_{m,j}))\right|
       \leq C_{h,c}\sqrt{\Delta_{m,j}}.                  \tag{9}
\]
This constant is independent of both $j$ and $m$.
At fixed $j$ take the subnet limit, using (M):
\[
 |\nu(h)-M(w_j)(h)|
       \leq C_{h,c}\sqrt{E(w_j)-E(w)}.
\]
Then take $j$ to its endpoint. Canonical weak convergence in (M)
gives $\nu(h)=M(w)(h)$. Model functions are uniformly dense in
$C(K)$, and masses agree, so this identifies $\nu$.
Compactness of bounded-mass Radon measures and uniqueness of all
clusters prove the whole weak limit (6).
\end{proof}

\section{Finite nef locality and bounded original locality}
Finite bounded relatively nef models on an actual nef $A$ satisfy
strict-winning locality. To see this within the exact G0 grant,
choose an actual ample $H$ on the generic carrier and add the same
zero metric in class $\epsilon H$ to both scalar models and their
maximum. Their actual generic bundle $A+\epsilon H$ is ample.
G0 finite ample locality applies; the winning scalar inequality
is unchanged. On each fixed common determination the raw measures
are finite polynomials in $\epsilon$, with every central multiplicity.
Let rational $\epsilon\downarrow0$. Positivity and the finite
coefficient formulas give the asserted finite nef identity.
This passage imports no unbounded ample estimate.

For $a,b,w$ of the previous section and a continuous original test
$h$ with compact support in $\{a>b\}\cap U$, compactness gives a
positive minimum for $a-b$ on that support.
The moving quotient gap converges uniformly there, so
$\psi_{m,f}>\psi_{m,g}$ at every pushed atom where $h$ is nonzero,
eventually. Finite nef locality gives equality of the $h$ integrals
of $\MA(\rho_m)$ and $\MA(\psi_{m,f})$.
Pass both weak limits using (6) and (M).
Radon uniqueness on the open subset of $U$ yields
\[
 \mathbf1_{\{a>b\}\cap U}M(\max(a,b))
       =\mathbf1_{\{a>b\}\cap U}M(a).                   \tag{10}
\]
The analogous identity holds where $b>a$.
The bounded winner is open in $U$ because these model envelopes are
continuous there. The unbounded extension below will not use
openness of its winner.

\section{All decreasing bases and uniformly small ordered errors}
\begin{lemma}[Decreasing weak continuity]\label{decreasing}
Under (S)--(M)--(A), every genuine decreasing net
$u_i\downarrow u\in\EE$ has $M(u_i)\to M(u)$ weakly.
\end{lemma}
\begin{proof}
For fixed $h\in C(K)$ and fixed $t$, the envelopes
$P_\theta(u_i+th)$ decrease to $P_\theta(u+th)$.
Their decreasing endpoint is genuine by (S), is below $u+th$
at every original point, and is above its envelope; hence it equals
that envelope. Its generic endpoint is finite because it lies
between $u\pm|t|\|h\|_\infty$.
Their energies converge decreasingly.
The concave functions $H_i(t)=E(P_\theta(u_i+th))$ thus converge
at each fixed $t$ to $H(t)$.
Their derivatives at zero are $\int h\,dM(u_i)$ by (A).
For $t>0$, the two secants at $-t,0,t$ bracket each such derivative.
First take the net limit at this fixed $t$, then let $t\downarrow0$.
Both derivative signs of $H$ agree by (A), so the bracket shrinks
to $\int h\,dM(u)$. This is weak convergence on the original $K$.
\end{proof}

For any canonical majorants $a_j\downarrow u$, monotonicity gives
\[
 0\leq\delta(a_j,u)\leq\frac{E(a_j)-E(u)}2.
\]
Consequently (1) gives numbers $\eta_j(D)\downarrow0$ such that
\[
 \sup_{B(z)\leq D}\int_U(a_j-u)\,dM(z)
       \leq\eta_j(D).                                  \tag{11}
\]
For example use (1) with $B(a_j)\leq B(u)$ and the displayed
energy gap. This is a uniform estimate for the \emph{fixed} test
majorant error against every output of cost at most $D$.
It follows from the absolute modulus, not from weak convergence
of any indicator.

\section{Strict Borel locality for all genuine pairs}
\begin{theorem}[Strict Borel locality]\label{borel}
Under the explicit inputs (S)--(M)--(A), every normalized
$u,v\in\EE$ satisfies
\[
 \mathbf1_{\{u>v\}}M(\max(u,v))
       =\mathbf1_{\{u>v\}}M(u),\qquad
 \mathbf1_{\{v>u\}}M(\max(u,v))
       =\mathbf1_{\{v>u\}}M(v).                         \tag{12}
\]
These are identities of original Borel measures.
\end{theorem}
\begin{proof}
Take canonical model envelopes $u_i\downarrow u$, $v_i\downarrow v$
on a product-directed index and put $w_i=\max(u_i,v_i)$,
$w=\max(u,v)$. Then $w_i\downarrow w$ genuinely and all output
costs are bounded by one $D$.
Lemma~\ref{decreasing} passes $M(w_i)$ and $M(u_i)$ weakly to
$M(w)$ and $M(u)$.

Let $\beta:\mathbb R\to[0,1]$ be zero on $(-\infty,1]$,
one on $[2,\infty)$ and linear between; its Lipschitz constant is one.
Fix $h\in C(K)$ with compact support in $U$.
For a fixed model index $j$ and $\varepsilon>0$, define on $U$
\[
 b_{j,\varepsilon}=\beta((u_j-v_j)/\varepsilon).
\]
The fixed function $h b_{j,\varepsilon}$ extends continuously by zero
to $K$. If $i\geq j$, $b_{j,\varepsilon}>0$ and $u_i\leq v_i$,
then $u_j-v_j>\varepsilon$, $v_i\leq v_j$, and hence
$u_j-u_i>\varepsilon$.
On the region $u_i>v_i$, identity (10) cancels the two measures.
On its complement, the preceding inequality and (11) bound both
positive measures. Thus
\[
 \left|\int h b_{j,\varepsilon}\,d(M(w_i)-M(u_i))\right|
       \leq \frac{2\|h\|_\infty}{\varepsilon}\eta_j(D).
                                                               \tag{13}
\]
Here $u_j-u_i\leq u_j-u$; no moving set is used as a weak test.
Take $i$ to its endpoint at this fixed $j,\varepsilon$.
Weak convergence passes the fixed continuous test in (13).

At the two now fixed measures $M(w),M(u)$ define
$b_\varepsilon=\beta((u-v)/\varepsilon)$ where both relative
costs are finite. Their complements have zero mass by (A).
The Lipschitz bound gives
\[
 |b_{j,\varepsilon}-b_\varepsilon|
 \leq \varepsilon^{-1}\big((u_j-u)+(v_j-v)\big)
\]
on this full-measure set. Estimate (11), separately for the
$u$ and $v$ nets, proves absolute $L^1$ convergence at these
two fixed measures. Sending $j$ to its endpoint in (13) gives
$\int h b_\varepsilon\,d(M(w)-M(u))=0$.

Finally take an ordinary positive sequence $\varepsilon\downarrow0$.
Bounded dominated convergence at the two fixed measures gives
$b_\varepsilon\to\mathbf1_{\{u>v\}}$ almost everywhere.
Thus their restrictions to that original Borel set have equal
integrals against every compactly supported continuous test in $U$.
Radon uniqueness gives (12) on $U$; all three measures miss $Z$
by (M), so it holds on $K$. Interchange $u,v$ for the other identity.
\end{proof}

\begin{theorem}[Strict Borel comparison]
For every genuine $u,v\in\EE$, under the same conditional inputs,
\[
 \int_{\{u<v\}}dM(v)\leq\int_{\{u<v\}}dM(u).             \tag{14}
\]
\end{theorem}
\begin{proof}
For $\varepsilon>0$ let $w_\varepsilon=\max(u,v-\varepsilon)$.
Finite energy and genuine membership follow by max closure and
energy monotonicity; restore constants when needed.
By (12) and constant invariance, its measure agrees with $M(v)$
on $\{u<v-\varepsilon\}$ and with $M(u)$ on
$\{u>v-\varepsilon\}$. All raw masses are $V$. Hence
\[
 M(v)\{u<v-\varepsilon\}
 \leq V-M(u)\{u>v-\varepsilon\}
 =M(u)\{u\leq v-\varepsilon\}.
\]
Send an ordinary sequence $\varepsilon\downarrow0$.
Both sets increase to $\{u<v\}$ off the null set where a relative
cost is infinite. Finite-measure continuity from below proves (14).
No winner is assumed open, and no indicator passes through weak convergence.
\end{proof}

\section{Every genuine plurisubharmonic pole is null}
\begin{proposition}[A finite-energy pole covering]\label{polar}
Under (S)--(C), for every normalized genuine $q\in\operatorname{PSH}(\theta)$,
there is an actual $p\in\EE$ whose pole set contains
$\{q=-\infty,\ R>-\infty\}$.
Consequently every output in (A) and every intrinsic target gives
$\{q=-\infty\}$ zero mass.
\end{proposition}
\begin{proof}
Put $N_j=16^j$, $r_j=8^{-j}$ and $\lambda_j=2^{-j}$, $j\geq1$.
The genuine truncation $q_j=\max(q,R-N_j)$ lies in $\EE$ since
$R-N_j\leq q_j\leq R$, and $B(q_j)\leq VN_j$.
Set $a_j=(1-r_j)R+r_jq_j$.
Contraction gives $B(a_j)\leq C_nV4^{-j}$.
The genuine finite sums
\[
 p_m=\sum_{j=1}^m\lambda_j a_j
           +\left(1-\sum_{j=1}^m\lambda_j\right)R
\]
decrease all-point, have generic value zero, and by concavity satisfy
$E(p_m)\geq-C_nV\sum_{j\geq1}8^{-j}$.
Decreasing closure and energy continuity in (S) give a genuine
finite-energy endpoint $p$.
Where $q=-\infty$ and $R$ is finite, $R-q_j=N_j$ and hence
$R-p_m=\sum_{j=1}^m\lambda_j r_jN_j=m$.
Thus $p=-\infty$ there.
Integrability of this one actual relative finite-energy test forces
zero measure on that set for every stated measure.
The remaining set $\{R=-\infty\}$ is already null by (M), or by
the intrinsic-target definition. This proves the assertion.
\end{proof}

\section{Limits of the result and source application scope}
All positive-volume conclusions here remain \emph{conditional}.
The whole scalar and moving foundation package has not been accepted.
The absolute midpoint modulus and derivative passage are unreviewed
conditional results of the whole round159 and continuation200 bodies.
This manuscript proves the bounded moving-max energy and measure
identification and its Borel-selection argument under those explicit inputs.
It does not turn those inputs into proved foundations.

The pole result concerns genuine $\theta$-plurisubharmonic pole sets.
It does not assert that every raw-supremum negligible discrepancy is
such a pole set. A proof of the requisite negligible-set theorem remains
separate. All unbounded mixed products, simultaneously changing physical
weighted approximations, their independence and their equality with this
outer map remain separate original obligations. No additional zero-volume
range is established by this positive-volume paper.
Original P16.5 therefore remains open; this draft has no independent
acceptance and causes no change to the solved-original count.

G0's complete supplied actual ample locality unit, inclusive saved
LF260--335, was read with its retained finite component coefficient
formula and complete locality/comparison statement. Only its bounded
finite ample assertion is applied here; a finite polynomial limit
handles nef models. The fixed-nef positive-product comparisons use the
ordinary algebraic BFJ theory with its characteristic-zero hypotheses,
not an unbounded big-class NA equation.
The complete saved BFJ v2 positive-product unit, inclusive LF1--132,
and volume/differentiation unit LF1--101 were read for this application.
Their original renderer coordinates are 310--441 and 445--545.
They include the nef agreement and characterization by dominated nef
competitors used in (2)--(4); no multilinearity of big positive
products is asserted.
The trace inequality (5), uniform base-twist generation, ordered
finite gap (8), fixed ramp estimate (13), and convex pole series are
proved above. The related whole unreviewed manuscript bodies have
been read; their source ledgers are provenance, not primary reading.
No new original PDF-byte hash, whole-primary reading, or recursive
background proof reading is claimed.

\begin{thebibliography}{9}
\bibitem{BJ} S. Boucksom and M. Jonsson,
\emph{Global pluripotential theory over a trivially valued field},
\href{https://arxiv.org/pdf/1801.08229v4}{arXiv:1801.08229v4}.
\bibitem{BFJ} S. Boucksom, C. Favre and M. Jonsson,
\emph{Differentiability of volumes of divisors and a problem of Teissier},
\href{https://arxiv.org/html/math/0608260v2}{arXiv:math/0608260v2}.
\end{thebibliography}
\end{document}
